\uuid{obQG}
\titre{Équation droite/plan}
\auteur{bodin ; hanani ; mzari}
\organisation{exo7}
\licence{CC BY-NC-SA 4.0}
\datecreate{2020-01-20}
\niveau{L1}
\module{Géométrie}
\chapitre{Géométrie de l'espace}
\sousChapitre{Plans}
\difficulte{1}
\theme{Plans, équation droite, plan}
\exo7qcmid{lille-1.geometrie-espace.005}

\contenu{
\texte{
Pour ces questions, l'espace est muni d'un repère orthonormé direct  $(O, \vec{i}, \vec{j}, \vec{k})$.
}

\begin{question}
\qclassification{141.03}
\qtags{motcle=équation droite/plan}
Soit $P$ le plan passant par $A(-1,1,1)$ et dirigé par les vecteurs
 $\vec{u}(0,1,1)$ et $\vec{v}(1,0,1)$. Quelles sont les assertions vraies ?
\begin{answers}

    \good {Une représentation paramétrique de $P$ est :
 $$\left\{\begin{array}{ccl}x&=&-1+s\\y&=&1+t\\ z&=&1+t+s,\quad (t,s\in \Rr) \end{array}\right.$$}

    \bad {Une représentation paramétrique de $P$ est :
 $$\left\{\begin{array}{ccl}x&=&-1+t\\y&=&1+s\\ z&=&-1+t+s,\quad (t,s\in \Rr)\end{array}\right.$$}

   \bad{Une équation cartésienne de $P$ est $x+y+z=-1$.}

    \good{Une équation cartésienne de $P$ est $x+y-z=-1$.}

\end{answers}
\begin{explanations}
On peut trouver une équation cartésienne, à partir d'une représentation paramétrique,  en éliminant les paramètres.
\end{explanations}
\end{question}
}
